% ============================================================
%  MTH636 – Game Theory | Lecture 14
%  Domination by Mixed Strategies & Nash's Existence Theorem
% ============================================================
\documentclass[12pt, a4paper]{article}

% ============================================================
%  MTH636 -- Game Theory
%  Notes Preamble — Color-Coded, TikZ-Enhanced
% ============================================================

% --- Core packages ---
\usepackage[T1]{fontenc}
\usepackage[utf8]{inputenc}
\usepackage{lmodern}
\usepackage[margin=2.5cm, top=2.8cm, bottom=2.8cm]{geometry}
\usepackage{microtype}

% --- Math ---
\usepackage{amsmath, amssymb, amsthm, mathtools}
\usepackage{bm}        % bold math
\usepackage{cancel}    % strike-through in math

% --- Color & Boxes ---
\usepackage[dvipsnames, svgnames, x11names, table]{xcolor}
\usepackage[most]{tcolorbox}
\tcbuselibrary{breakable, skins, theorems}

% --- TikZ & Diagrams ---
\usepackage{tikz}
\usepackage{tikz-qtree}
\usetikzlibrary{arrows.meta, positioning, shapes.geometric, fit, backgrounds,
                decorations.pathreplacing, calc, matrix, trees}
\usepackage{pgfplots}
\pgfplotsset{compat=1.18}

% --- Tables ---
\usepackage{booktabs}
\usepackage{multirow}
\usepackage{array}
\usepackage{tabularx}
\usepackage{makecell}

% --- Typography & Layout ---
\usepackage{enumitem}
\usepackage{parskip}
\usepackage{fancyhdr}
\usepackage{titlesec}
\usepackage{titletoc}
\usepackage{hyperref}
\usepackage{graphicx}
\usepackage{float}
\usepackage{caption}

% --- Color Palette ---
\definecolor{defblue}{HTML}{1A5276}      % Definition header
\definecolor{defbluebg}{HTML}{D6EAF8}    % Definition background
\definecolor{thmgreen}{HTML}{1E8449}     % Theorem header
\definecolor{thmbg}{HTML}{D5F5E3}        % Theorem background
\definecolor{exorange}{HTML}{B7770D}     % Example header
\definecolor{exbg}{HTML}{FEF9E7}         % Example background
\definecolor{warnred}{HTML}{922B21}      % Warning/Key Insight header
\definecolor{warnbg}{HTML}{FADBD8}       % Warning background
\definecolor{noticegray}{HTML}{2C3E50}   % Note header
\definecolor{noticebg}{HTML}{EBF5FB}     % Note background
\definecolor{taketawayyellow}{HTML}{7D6608} % Key Takeaway header
\definecolor{takeawaybg}{HTML}{FDFEFE}   % Takeaway background
\definecolor{accentpurple}{HTML}{6C3483} % Accent / proof
\definecolor{lightgray}{HTML}{F2F3F4}    % shading

% ============================================================
%  TCOLORBOX ENVIRONMENTS
% ============================================================

% -- Definition (Blue) --
\newtcolorbox{defbox}[2][]{
  enhanced, breakable,
  colback=defbluebg, colframe=defblue, coltitle=white,
  fonttitle=\bfseries\sffamily,
  title={Definition: #2},
  attach boxed title to top left={yshift=-2mm, xshift=4mm},
  boxed title style={colback=defblue, rounded corners},
  arc=4pt, boxrule=1.2pt,
  left=6pt, right=6pt, top=8pt, bottom=6pt,
  #1
}

% -- Theorem / Proposition (Green) --
\newtcolorbox{thmbox}[2][]{
  enhanced, breakable,
  colback=thmbg, colframe=thmgreen, coltitle=white,
  fonttitle=\bfseries\sffamily,
  title={Theorem: #2},
  attach boxed title to top left={yshift=-2mm, xshift=4mm},
  boxed title style={colback=thmgreen, rounded corners},
  arc=4pt, boxrule=1.2pt,
  left=6pt, right=6pt, top=8pt, bottom=6pt,
  #1
}

% -- Proposition (Green variant) --
\newtcolorbox{propbox}[2][]{
  enhanced, breakable,
  colback=thmbg, colframe=thmgreen!70!black, coltitle=white,
  fonttitle=\bfseries\sffamily,
  title={Proposition: #2},
  attach boxed title to top left={yshift=-2mm, xshift=4mm},
  boxed title style={colback=thmgreen!70!black, rounded corners},
  arc=4pt, boxrule=1.2pt,
  left=6pt, right=6pt, top=8pt, bottom=6pt,
  #1
}

% -- Example (Orange/Yellow) --
\newtcolorbox{exbox}[2][]{
  enhanced, breakable,
  colback=exbg, colframe=exorange, coltitle=white,
  fonttitle=\bfseries\sffamily,
  title={Example: #2},
  attach boxed title to top left={yshift=-2mm, xshift=4mm},
  boxed title style={colback=exorange, rounded corners},
  arc=4pt, boxrule=1.2pt,
  left=6pt, right=6pt, top=8pt, bottom=6pt,
  #1
}

% -- Key Insight / Warning (Red) --
\newtcolorbox{warnbox}[2][]{
  enhanced, breakable,
  colback=warnbg, colframe=warnred, coltitle=white,
  fonttitle=\bfseries\sffamily,
  title={Key Insight: #2},
  attach boxed title to top left={yshift=-2mm, xshift=4mm},
  boxed title style={colback=warnred, rounded corners},
  arc=4pt, boxrule=1.2pt,
  left=6pt, right=6pt, top=8pt, bottom=6pt,
  #1
}

% -- Note (Gray-Blue) --
\newtcolorbox{notebox}[1][]{
  enhanced, breakable,
  colback=noticebg, colframe=noticegray!60, coltitle=white,
  fonttitle=\bfseries\sffamily,
  title={Note},
  attach boxed title to top left={yshift=-2mm, xshift=4mm},
  boxed title style={colback=noticegray!60, rounded corners},
  arc=4pt, boxrule=0.8pt,
  left=6pt, right=6pt, top=8pt, bottom=6pt,
  #1
}

% -- Key Takeaway --
\newtcolorbox{takebox}[1][]{
  enhanced, breakable,
  colback=takeawaybg, colframe=taketawayyellow, coltitle=white,
  fonttitle=\bfseries\sffamily,
  title={Key Takeaway},
  attach boxed title to top left={yshift=-2mm, xshift=4mm},
  boxed title style={colback=taketawayyellow, rounded corners},
  arc=4pt, boxrule=1.2pt,
  left=6pt, right=6pt, top=8pt, bottom=6pt,
  #1
}

% -- Intuition --
\newtcolorbox{intubox}[1][]{
  enhanced, breakable,
  colback=accentpurple!8, colframe=accentpurple!60, coltitle=white,
  fonttitle=\bfseries\sffamily,
  title={Intuition},
  attach boxed title to top left={yshift=-2mm, xshift=4mm},
  boxed title style={colback=accentpurple!80, rounded corners},
  arc=4pt, boxrule=0.8pt,
  left=6pt, right=6pt, top=8pt, bottom=6pt,
  #1
}

% ============================================================
%  PAYOFF MATRIX HELPER
%  Usage: inside tabular, use \payoff{a}{b} for (a,b)
% ============================================================
\newcommand{\payoff}[2]{$(#1,\;#2)$}
\newcommand{\nash}[2]{\underline{$#1$},\underline{$#2$}}  % underline NE payoffs
\newcommand{\Nset}{\mathbb{N}}
\newcommand{\Rset}{\mathbb{R}}
\newcommand{\Sset}{\mathcal{S}}
\newcommand{\Ui}{u_i}

% ============================================================
%  SECTION FORMATTING
% ============================================================
\titleformat{\section}
  {\large\bfseries\sffamily\color{defblue}}
  {\thesection}{1em}{}[\titlerule]

\titleformat{\subsection}
  {\normalsize\bfseries\sffamily\color{thmgreen!80!black}}
  {\thesubsection}{1em}{}

\titleformat{\subsubsection}
  {\normalsize\itshape\sffamily\color{exorange!80!black}}
  {\thesubsubsection}{1em}{}

% ============================================================
%  HEADER / FOOTER
% ============================================================
\pagestyle{fancy}
\fancyhf{}
\fancyhead[L]{\sffamily\small\color{defblue} MTH636 -- Game Theory}
\fancyhead[R]{\sffamily\small\color{defblue} \leftmark}
\fancyfoot[C]{\sffamily\small\thepage}
\renewcommand{\headrulewidth}{0.4pt}

% ============================================================
%  HYPERREF SETUP
% ============================================================
\hypersetup{
  colorlinks=true,
  linkcolor=defblue,
  urlcolor=thmgreen,
  citecolor=accentpurple,
  bookmarksnumbered=true,
  pdftitle={MTH636 Game Theory Notes},
  pdfauthor={MTH636 Student Notes}
}

% ============================================================
%  MISC HELPERS
% ============================================================
\newcommand{\BR}{\text{BR}}        % Best Response
\newcommand{\NE}{\text{NE}}        % Nash Equilibrium
\newcommand{\SPNE}{\text{SPNE}}    % Subgame Perfect NE
\newcommand{\Dom}{\text{dom}}      % Dominance
\newcommand{\mixs}[1]{\sigma_{#1}} % Mixed strategy
\newcommand{\pures}[1]{s_{#1}}     % Pure strategy


% ============================================================
\begin{document}

% Title
\begin{center}
  {\LARGE\bfseries\sffamily\color{defblue} MTH636 -- Game Theory}\\[6pt]
  {\Large\sffamily Lecture 14: Domination by Mixed Strategies \& Nash's Theorem}\\[4pt]
\end{center}
\vspace{0.5cm}
\hrule height 1.5pt
\vspace{0.3cm}

% ============================================================
\section{Domination by Mixed Strategies}

A pure strategy can be strictly dominated not only by another pure strategy, but also by a \textit{mixed} strategy.

\begin{exbox}{Strict Domination by a Mixed Strategy}
  Consider the following game matrix:
  \begin{center}
  \renewcommand{\arraystretch}{1.4}
  \begin{tabular}{c|ccc}
    $P_1 \setminus P_2$ & L & C & R \\
    \hline
    T & \payoff{6}{2} & \payoff{0}{6} & \payoff{4}{4} \\
    M & \payoff{2}{12} & \payoff{4}{3} & \payoff{2}{5} \\
    B & \payoff{0}{6} & \payoff{10}{0} & \payoff{2}{2} \\
  \end{tabular}
  \end{center}
  
  Observe that the pure strategy $M$ for Player 1 is not strictly dominated by either $T$ or $B$ alone. However, consider the mixed strategy $\sigma_1 = [\frac{1}{2}(T), 0(M), \frac{1}{2}(B)]$. Let's calculate the expected payoff of $\sigma_1$ against each of Player 2's pure strategies:
  \begin{itemize}
      \item Against L: $U_1(\sigma_1, L) = \frac{1}{2}(6) + \frac{1}{2}(0) = 3 > 2 = U_1(M, L)$
      \item Against C: $U_1(\sigma_1, C) = \frac{1}{2}(0) + \frac{1}{2}(10) = 5 > 4 = U_1(M, C)$
      \item Against R: $U_1(\sigma_1, R) = \frac{1}{2}(4) + \frac{1}{2}(2) = 3 > 2 = U_1(M, R)$
  \end{itemize}
  Since $U_1(\sigma_1, s_2) > U_1(M, s_2)$ for all $s_2 \in S_2$, we have:
  \[
    U_1(\sigma_1, \sigma_2) > U_1(M, \sigma_2) \quad \forall \sigma_2 \in \Sigma_2
  \]
  Thus, $M$ is strictly dominated by the mixed strategy $\sigma_1$.
\end{exbox}

\begin{defbox}{Strict Domination (General Definition)}
  A strategy $\sigma_i$ of player $i$ \textbf{strictly dominates} a pure strategy $s_i$ of player $i$ if:
  \[
    U_i(\sigma_i, \sigma_{-i}) > U_i(s_i, \sigma_{-i}) \quad \forall \sigma_{-i} \in \Sigma_{-i}
  \]
  Equivalently, it suffices to check against all pure strategies of the opponent:
  \[
    U_i(\sigma_i, s_{-i}) > U_i(s_i, s_{-i}) \quad \forall s_{-i} \in S_{-i}
  \]
\end{defbox}

\begin{thmbox}{Strictly Dominated Strategies in Nash Equilibrium}
  Suppose for player $i$, a strategy $\sigma_i$ strictly dominates a pure strategy $s_i$ of player $i$. Then $s_i$ \textbf{does not get a positive probability} in any Nash Equilibrium of the game.
\end{thmbox}

\begin{proof}
  Suppose $\overline{\sigma}$ is a mixed strategy profile and $\overline{\sigma}_i(s_i) > 0$. We will show that there exists a mixed strategy $\overline{\overline{\sigma}}_i$ such that $\overline{\overline{\sigma}}_i$ strictly dominates $\overline{\sigma}_i$ against any opponent strategy profile, meaning $\overline{\sigma}$ cannot be a best response and thus cannot be part of an NE.
  
  Construct the new strategy $\overline{\overline{\sigma}}_i$ by taking the probability mass assigned to the dominated strategy $s_i$ and reassigning it proportionally to the components of the dominating strategy $\sigma_i$. Specifically, for any pure strategy $t_i \in S_i$:
  \[
    \overline{\overline{\sigma}}_i(t_i) = \begin{cases}
      \overline{\sigma}_i(s_i) \cdot \sigma_i(s_i) & \text{if } t_i = s_i \\
      \overline{\sigma}_i(t_i) + \overline{\sigma}_i(s_i) \sigma_i(t_i) & \text{if } t_i \neq s_i
    \end{cases}
  \]
  By calculating the expected utility, we can sum the terms:
  \begin{align*}
    U_i(\overline{\overline{\sigma}}_i, \sigma_{-i}) &= \sum_{t_i \in S_i} \overline{\overline{\sigma}}_i(t_i) U_i(t_i, \sigma_{-i}) \\
    &= \sum_{t_i \neq s_i} \left[ \overline{\sigma}_i(t_i) + \overline{\sigma}_i(s_i)\sigma_i(t_i) \right] U_i(t_i, \sigma_{-i}) + \overline{\sigma}_i(s_i)\sigma_i(s_i)U_i(s_i, \sigma_{-i}) \\
    &= \sum_{t_i \neq s_i} \overline{\sigma}_i(t_i) U_i(t_i, \sigma_{-i}) + \overline{\sigma}_i(s_i) \left( \sum_{t_i \neq s_i} \sigma_i(t_i) U_i(t_i, \sigma_{-i}) + \sigma_i(s_i)U_i(s_i, \sigma_{-i}) \right) \\
    &= \sum_{t_i \neq s_i} \overline{\sigma}_i(t_i) U_i(t_i, \sigma_{-i}) + \overline{\sigma}_i(s_i) U_i(\sigma_i, \sigma_{-i})
  \end{align*}
  Since $\sigma_i$ strictly dominates $s_i$, we have $U_i(\sigma_i, \sigma_{-i}) > U_i(s_i, \sigma_{-i})$. Substituting this inequality gives:
  \begin{align*}
    U_i(\overline{\overline{\sigma}}_i, \sigma_{-i}) &> \sum_{t_i \neq s_i} \overline{\sigma}_i(t_i) U_i(t_i, \sigma_{-i}) + \overline{\sigma}_i(s_i) U_i(s_i, \sigma_{-i}) \\
    &= U_i(\overline{\sigma}_i, \sigma_{-i})
  \end{align*}
  Thus, $\overline{\overline{\sigma}}_i$ yields a strictly higher payoff than $\overline{\sigma}_i$, meaning $\overline{\sigma}_i$ cannot be a best response, and thus cannot be part of an NE.
\end{proof}

% ============================================================
\section{Proof of Nash's Theorem}

\begin{thmbox}{Nash's Existence Theorem}
  Suppose $G = \langle N, (S_i)_{i \in N}, (u_i)_{i \in N} \rangle$ is a strategic form game where all the players have finitely many pure strategies. Then $G$ has at least one \textbf{mixed strategy Nash Equilibrium}.
\end{thmbox}

To prove this, we rely on Brouwer's Fixed Point Theorem.

\begin{thmbox}{Brouwer's Fixed Point Theorem}
  Let $X$ be a convex and compact subset of $\mathbb{R}^d$ for some $d \in \mathbb{N}$, and let $f: X \to X$ be a continuous function. Then $f$ has a fixed point, i.e., $\exists x \in X$ such that $f(x) = x$.
\end{thmbox}

\textbf{Proof of Nash's Theorem:} \\
Let $m_i = |S_i|$. The space of mixed strategies for player $i$ is the simplex:
\[
  \Sigma_i = \left\{ (x_1, \dots, x_{m_i}) \mid x_j \geq 0 \ \forall j, \ \sum_{j=1}^{m_i} x_j = 1 \right\}
\]
\textbf{Claim 1:} $\Sigma_i$ is a convex and compact subset of $\mathbb{R}^{m_i}$. \\
\textbf{Claim 2:} Suppose $A$ and $B$ are convex and compact subsets of $\mathbb{R}^{d_1}$ and $\mathbb{R}^{d_2}$ respectively. Then $A \times B$ is a convex and compact subset of $\mathbb{R}^{d_1 + d_2}$. \\
From Claims 1 and 2, the set of all mixed strategy profiles $\Sigma = \Sigma_1 \times \Sigma_2 \times \dots \times \Sigma_n$ is a convex and compact subset of $\mathbb{R}^{\sum m_i}$.

We will construct a function $f: \Sigma \to \Sigma$ such that:
\begin{enumerate}
    \item[(i)] $f$ is continuous.
    \item[(ii)] Every fixed point of $f$ is a Nash Equilibrium of $G$.
\end{enumerate}
By Brouwer's Fixed Point Theorem, such a function must have a fixed point, which will imply the existence of a Nash Equilibrium.

For a strategy profile $\sigma \in \Sigma$, let $\delta_i^j(\sigma)$ measure the ``gain'' a player gets by deviating to their $j$-th pure strategy $s_i^j$:
\[
  \delta_i^j(\sigma) = \max \left\{ 0, \ U_i(s_i^j, \sigma_{-i}) - U_i(\sigma) \right\}
\]
\textbf{Claim 3:} $\delta_i^j$ is a continuous function of $\sigma$. \\
\textit{Proof of Claim 3:} The expected utility $U_i(\sigma)$ is a multilinear function of the probabilities in $\sigma$, hence it is continuous. $U_i(s_i^j, \sigma_{-i})$ is similarly multilinear in $\sigma_{-i}$ and thus continuous in $\sigma$. Since the difference of continuous functions is continuous, and the $\max(0, \cdot)$ function is continuous, $\delta_i^j(\sigma)$ is continuous.

(The full construction of $f$ continues by using these $\delta_i^j$ values to strictly increase the probabilities of pure strategies that yield better-than-average payoffs, normalized so they sum to 1).

\end{document}
